%% kapa: exam=1 round=27 year=2016 period=1 form=A subject=civil_law subject_ko=민법 seq=25 number=25
%% source: Q-Net 감정평가사 기출문제 (cst00309) artlSeq=5169402 file=extracted/1008/5169402/제27회 감정평가사 1차 문제/제27회 감정평가사 1교시문제(A형).pdf
%% conversion: kordoc 4.24.51 -> tools/gen1.py (auto); review=proofread-sonnet; flags=none
\begin{question}{25}
중간생략등기에 관한 설명으로 옳지 않은 것은? (다툼이 있으면 판례에 따름)
\begin{choices}
  \item 甲이 신축한 건물을 乙이 매수한 후, 당사자들의 합의에 따라 경료된 乙 명의의 보존등기는 유효하다.
  \item 토지거래허가구역 내 토지가 甲에서 乙, 乙에서 丙으로 매도되고 중간생략등기의 합의가 있더라도, 丙이 자신과 甲을 매매 당사자로 하는 토지거래허가를 받아 丙 앞으로 경료된 소유권이전등기는 무효이다.
  \item 매도인 甲, 중간매수인 乙, 최후매수인 丙 사이에 중간생략등기에 대한 전원의 합의가 없는 경우, 丙은 甲에 대하여 직접 자기에게 이전등기를 청구할 수 없다.
  \item 매도인 甲, 중간매수인 乙, 최후매수인 丙이 甲으로부터 丙으로 이전등기를 해주기로 전원 합의한 경우, 乙이 대금을 지급하지 않더라도 甲은 丙에게 소유권이전등기를 해주어야 한다.
  \item 매도인 甲, 중간매수인 乙, 최후매수인 丙이 甲으로부터 丙으로 이전등기를 해주기로 전원 합의한 경우에도 乙은 甲에 대한 등기청구권을 잃지 않는다.
\end{choices}
\end{question}
\answer{4}
